\documentclass[10pt,a4paper]{amsart}
\usepackage[margin=19mm]{geometry}
\usepackage{amsmath,amssymb,mathtools}
\usepackage[scr=rsfs,cal=euler]{mathalfa}
\usepackage{xfrac}
\usepackage[unicode,hidelinks,bookmarks=false]{hyperref}
\newcommand{\ie}{i.e.\ }
\newcommand{\cf}{cf.\ }
\newcommand{\resp}{respectively\ }
\newcommand{\Subsection}{\subsection}
\providecommand{\nobreakdash}{\nobreak\hbox{-}}
\setlength{\emergencystretch}{2em}
\multlinegap0pt
\usepackage{lmodern}
\usepackage{mathtools}
\usepackage{booktabs}
\usepackage{array}
\usepackage{enumitem}
\allowdisplaybreaks
\newtheorem{theorem}{Theorem}[section]
\newtheorem{proposition}[theorem]{Proposition}
\newtheorem{lemma}[theorem]{Lemma}
\newtheorem{corollary}[theorem]{Corollary}
\newtheorem{criterion}[theorem]{Criterion}
\newtheorem{definition}[theorem]{Definition}
\newtheorem{question}[theorem]{Question}
\newtheorem{example}[theorem]{Example}
\newtheorem{remark}[theorem]{Remark}
\DeclareMathOperator{\depth}{depth}
\DeclareMathOperator{\Ass}{Ass}
\DeclareMathOperator{\ann}{ann}
\DeclareMathOperator{\rank}{rank}
\DeclareMathOperator{\Span}{span}
\DeclareMathOperator{\Soc}{Soc}
\DeclareMathOperator{\Gal}{Gal}
\DeclareMathOperator{\Aut}{Aut}
\DeclareMathOperator{\Hom}{Hom}
\DeclareMathOperator{\Ext}{Ext}
\DeclareMathOperator{\res}{res}
\DeclareMathOperator{\Nil}{Nil}
\DeclareMathOperator{\Spec}{Spec}
\DeclareMathOperator{\Supp}{Supp}
\DeclareMathOperator{\Proj}{Proj}
\DeclareMathOperator{\im}{im}
\DeclareMathOperator{\kerop}{ker}
\newcommand{\F}{\mathbb F}
\newcommand{\kbar}{\overline{\mathbb F}_{2}}
\newcommand{\C}{\mathbb C}
\newcommand{\Q}{\mathbb Q}
\newcommand{\Z}{\mathbb Z}
\newcommand{\SG}[2]{\operatorname{SmallGroup}(#1,#2)}
\newcommand{\EA}{\mathcal E}
\newcommand{\Nrad}{\sqrt{0}}
\newcommand{\A}{A}
\newcommand{\Azer}{A_0}
\newcommand{\Res}{\operatorname{Res}}
\begin{document}
\title[Centralizer Excess as an Obstruction to Carlson's Depth Conj]{Centralizer Excess as an Obstruction to Carlson's Depth Conjecture}
\author{Secondary 13C15, 20D15, 13D45, 18G50; Xinan Dai, Kuok Fai Chao}
\date{}
\maketitle
\noindent\emph{Source and licence.} Centralizer Excess as an Obstruction to Carlson's Depth Conjecture. Version arXiv:2607.27672v3 (CC BY 4.0). This material is distributed under the Creative Commons Attribution 4.0 International licence, http://creativecommons.org/licenses/by/4.0/, as recorded in the arXiv OAI licence metadata for this exact version and shown on the abstract page https://arxiv.org/abs/2607.27672v3. Full rights notice: licenses/arxiv-2607.27672-CC-BY-4.0.txt.\par\smallskip
\noindent\textit{Editorial scope.} Counted unit: the original \texttt{theorem} environment \texttt{thm:semidirect} at source lines 719--752 together with its complete proof at lines 754--781; the delivered document keeps source lines 708--782 so that every referenced label and every citation resolves inside it. Native editable source is the paper's own concatenated TeX; the AMS wrapper is editorial formatting only. No publisher class, upstream script or unrelated artwork is redistributed.
\section{The finite group and its elementary-abelian geometry}
\label{sec:group}

The Small Groups identifier and the rank-two centralizer table suffice for the counterexample.  We record a split extension because it gives convenient coordinates for the maximal elementary abelian subgroups.

Let \(F\) be the free group on \(x_1,\dots,x_7\), and let \(G=F/N\) be the presentation reproduced in \cite[Appendix~A.1]{DaiEtAl2026}.  Set
\[
u=x_2x_5,\qquad v=x_4,
\qquad s=x_1,\qquad t=x_3.
\]


\begin{theorem}[Explicit semidirect decomposition]
\label{thm:semidirect}
Let
\[
N_0=\langle u,v\rangle,
\qquad
H_0=\langle s,t\rangle.
\]
Then
\[
N_0\cong C_4\times C_4,
\qquad
H_0\cong C_4\times C_2,
\]
\(N_0\trianglelefteq G\), \(N_0\cap H_0=1\), and \(G=N_0H_0\).  Hence
\[
G\cong(C_4\times C_4)\rtimes_{\rho}(C_4\times C_2).
\]
With respect to the ordered basis \((u,v)\) of \((\Z/4\Z)^2\), and the convention \(n^h=h^{-1}nh\), the action is
\[
\rho(s)=
\begin{pmatrix}
1&2\\
3&3
\end{pmatrix},
\qquad
\rho(t)=
\begin{pmatrix}
1&0\\
2&1
\end{pmatrix}
\quad\text{in }\operatorname{GL}_2(\Z/4\Z).
\]
\end{theorem}

\begin{proof}
Collection in the finite presentation gives
\[
|u|=|v|=4,
\qquad uv=vu,
\qquad \langle u\rangle\cap\langle v\rangle=1,
\]
so \(|N_0|=16\) and \(N_0\cong C_4\times C_4\).  The relations
\[
s^2=x_5,\qquad x_5^2=1,
\qquad t^2=1,
\qquad [s,t]=1
\]
show that \(H_0\cong C_4\times C_2\) and \(|H_0|=8\).  Exact enumeration of the presented group gives \(|G|=128\), while
\[
N_0\cap H_0=1.
\]
The conjugation relations are
\[
u^s=uv^{-1},\qquad v^s=u^2v^{-1},
\]
and
\[
u^t=uv^2,\qquad v^t=v.
\]
Thus \(H_0\) normalizes \(N_0\), and the displayed matrices are obtained by reading the exponent vectors of these images as columns.  Since
\(|N_0H_0|=|N_0||H_0|=128\), one has \(G=N_0H_0\).  This proves the split semidirect decomposition.
\end{proof}


\begin{thebibliography}{99}
\bibitem[DaiEtAl2026]{DaiEtAl2026} X.~Dai, W.~Deng, Y.~Shi, T.~Wu, and Y.~Yang, An exact counterexample to Carlson's associated-prime depth conjecture from a group of order 128, arXiv preprint arXiv:2607.23732, 2026.
\end{thebibliography}
\end{document}
